\section{Conclusion}
\label{sec:conclusion}

We have introduced and validated the resolution selection rule for state
legislative redistricting: use block-group resolution when
$k/n_{\rm tracts} > 0.05$. The rule is derived from the requirement that
each district contain at least 20 base geographic units on average, ensuring
population balance feasibility, compactness optimisation richness, and contiguity
tractability.

The rule is empirically validated on three state house chambers spanning the
relevant range of $k/n$ values. For Washington's House ($k/n = 0.067$, above
threshold), block-group resolution improves mean PP from $0.368$ to $0.398$
and reduces maximum population deviation from $0.80\%$ to $0.40\%$.
For Texas ($k/n = 0.028$, below threshold), the improvement is marginal
($+0.008$ PP, $-0.06$\,pp deviation) and does not justify the $4\times$
runtime increase. For California ($k/n = 0.009$, far below threshold),
block-group resolution provides essentially no benefit.

The resolution recommendation table (Section~\ref{sec:table}) shows that
39 of 50 state lower chambers require block-group resolution, and 3 require
block resolution---a much higher proportion than might be expected given that
the congressional literature universally uses tract resolution.

\paragraph{MAUP implications.}
The MAUP analysis (Section~\ref{sec:maup}) shows that resolution choice
is not merely a technical decision for high-$k/n$ chambers: it can shift
partisan outcomes by $\pm 2$--$4$ seats. For algorithmically neutral
redistricting, the resolution selection rule must be stated and applied
consistently before any maps are drawn. The $k/n > 0.05$ rule provides
an objective, data-driven standard that is independent of any partisan
calculations.

\paragraph{Data infrastructure.}
The primary infrastructure requirement for the F track is the block-group
adjacency dataset. Building this dataset for all 50 states requires:
\begin{verbatim}
python scripts/data/geography/build_unit_adjacency.py \
  --resolution block_group --year 2020 --states all
\end{verbatim}
This is a one-time cost (approximately 4~GB, 2--4 hours) that enables all
future state legislative redistricting analyses at appropriate resolution.

\paragraph{Future work.}
The resolution selection rule should be extended to upper chambers (state senates)
and to bicameral nesting (F.2). Senate chambers have larger district sizes than
house chambers, so fewer states require block-group resolution for the senate.
The interaction between resolution selection and nesting (F.2) is a complex
constraint: nesting requires that house and senate maps use compatible base
units, which may require both chambers to use block-group resolution even if
only one chamber individually requires it.
